\documentclass[12pt]{article}
\usepackage[utf8]{inputenc}

\usepackage{../mystyle}

\title{Hilbert Scheme}
\author{Joseph Sullivan}
\date{September 2025}

\DeclareMathOperator{\Kom}{Kom}
\DeclareMathOperator{\Cyl}{Cyl}
\DeclareMathOperator{\Stab}{Stab}
\DeclareMathOperator{\Slice}{Slice}
\DeclareMathOperator{\chern}{ch}
\DeclareMathOperator{\todd}{td}
\DeclareMathOperator{\Spin}{Spin}

% \DeclareMathOperator{\Quot}{Quot}
\DeclareMathOperator{\Hilb}{Hilb}
\DeclareMathOperator{\Gr}{Gr}

\DeclareMathOperator{\ext}{ext}
%\DeclareMathOperator{\hom}{hom}

\newcommand{\cQ}{\mathcal{Q}}
\newcommand{\cZ}{\mathcal{Z}}
\newcommand{\LCom}{\mathcal{LC}}

\newcommand{\rddots}{\text{\reflectbox{$\ddots$}}}


\begin{document}

\maketitle
\tableofcontents

\section{Some Moduli Functors and Families}

\subsection{Technicalities}

We are going to be defining some moduli functors of interest, but in order to define our notions of families, we will need some technical definitions.

\begin{defn}
    A ring map $R\to S$ is of \textbf{finite presentation} if there exist finitely many $R$-algebra generators for $S$ with finitely many relations, so
    \[S\cong R[x_1,\dots,x_n]/(f_1,\dots,f_m)\]
    and the map from $R$ is the inclusion to the polynomial ring (and the projection).
\end{defn}

\begin{defn}
    Let $X\xrightarrow{f} S$.

    \begin{itemize}
        \item We say $f$ is of \textbf{finite presentation} at $x\in X$ if there exists an affine open neighborhood $\Spec(A) = U\subseteq X$ of $x$ and affine open $\Spec(R)=V\subseteq S$ with $f(U)\subseteq V$ such that the induced ring map $R\to A$ is of finite presentation.

        \item Then, $f$ is \textbf{locally of finite presentation} if it is of finite presentation at every point of $X$ (so this is local on the base).

        \item Finally, $f$ is \textbf{finite presentation} is is it locally of finite presentation, quasi-compact, and quasi-separated.
    \end{itemize}
\end{defn}

\begin{eg}
    When $S$ is locally Noetherian, locally of finite type already implies locally of finite presentation, since the kernel of the map on sections will be finitely generated by the Noetherian assumption.
\end{eg}

About morphisms of finite presentation, we have the following fact (which we will not prove).
\begin{thm}
    A morphism $f:X\to S$ is of finite presentation if and only if for all directed systems of affine schemes $T_i$,
    \[\operatorname{colim}\Hom_S(T_i,X) = \Hom_S(\lim T_i, X).\]
    (Note that we are taking in a colimit rather than a limit as is usual, and the limit ends up on the left argument, also unusual---this looks something like the definition of a compact object in a category).
\end{thm}

\begin{defn}
    Let $X\xrightarrow{f} S$, and $\cF$ a sheaf on $X$.
    
    We say $\cF$ is \textbf{flat} over $S$ at a point $x\in X$ if the stalk $\cF_x$ is a flat $\cO_{S,f(x)}$-module (by extension of scalars).

    We say $\cF$ is \textbf{flat} over $S$ if it is flat over $S$ at every point.

    We say $f$ is \textbf{flat} if $\cO_X$ is flat over $S$.
\end{defn}


\subsection{The Functors}

We first give a warm-up moduli functor.

\begin{defn}
    Let $V$ be a vector bundle of rank $n$ on a Noetherian scheme $S$, and $q$ an integer with $0<q<n$. Define the functor
    \begin{align*}
        \Gr(q,V): \Sch^\op &\to \Set\\
        (T\to S) &\mapsto \{\text{vector bundle quotients $V_T\twoheadrightarrow Q$ on $T$ of rank $q$}\}.
    \end{align*}
    We will call this the \textbf{Grassmannian}, especially after we establish that it is representable by a projective scheme over $S$. In the special case that $S=\Spec k$ (for $k$ a field or maybe $\Z$), then $V$ is a vector space (or an abelian group), and we recover the classical Grassmannian.
\end{defn}

\begin{defn}
    Let $S$ be a Noetherian scheme, $F$ a coherent sheaf on $\P_S^n$ that can be written as a quotient of $\cO_{\P^n_S}(-\ell)^{\oplus r}$ for some $l,r$. Finally, fix a polynomial $P\in \Q[z]$. Define the functor
    \begin{align*}
        \Quot^P(F/\P_S^n/S): (\Sch/S)^\op &\to \Set\\
        (T\to S) &\mapsto \left\{\parbox{22em}{quasi-coherent and finitely presented quotients\\ $F_T\twoheadrightarrow Q$ on $\P_T^n$ such that $Q$ is flat over $T$ and\\ $Q_t\defeq Q|_{\P^n_{\kappa(t)}}$ has Hilbert polynomial $P$ for all $t\in T$}\right\}
    \end{align*}
    which we will call the \textbf{Quot Scheme} (once we know it is represented by a scheme projective over $S$).
\end{defn}

In words, $\Quot^P(F/\P_S^n/S)$ is the moduli functor for quotients $Q$ of $F$ with a fixed Hilbert polynomial $P$. The condition that $F$ can be written as a quotient $\cO_{\P_S^n}(-\ell)^{\oplus r}$ is automatically satisfied for $S=\Spec A$ for a Noetherian ring $A$, by Serre vanishing (so this should be thought of as a technical condition necessary for the relative case).

We can then specialize to the setting where $F=\cO_{\P_S^n}$, so that quotients are exactly the (pushforwards of) structure sheaves of closed subschemes (i.e., those cut out by the kernel of the quotient map, the ideal sheaf). Then, $F_T$ is just $\cO_{\P^n_T}$, so $Q$ is a quotient $\cO_Z$ for a closed subscheme
\[Z\subseteq \P_T^n.\]
The condition that $\cO_Z$ is flat over $T$ is exactly what it means for the composition $Z\to \P_S^n \to S$ to be flat. Since $Z$ is a closed subscheme of projective space,
\[Z \to \P_S^n \to S\]
is separated (so in particular quasi-separated), and since $S$ is Noetherian, we similarly get $Z$ is quasi-compact. All this is to show that $Z\to S$ is finitely presented, where the finite presentation
\[\cO_{\P_T^n} \to \cO_Z\]
will give us a finitely type ideal sheaf, which will give us the finitely many algebra generators for the local model of the map $Z\to S$.

All this being said, we have the following definition.

\begin{defn}
    Let $S$ be a Noetherian scheme and $P\in \Q[z]$. Define the functor
    \begin{align*}
        \Hilb^P(\P_S^n/S): (\Sch/S)^\op &\to \Set\\
        (T\to S) &\mapsto \left\{\parbox{22em}{closed subschemes $Z\subseteq \P_T^n$ flat and finitely presented over $T$\\ such that $Z_t\subseteq \P_{\kappa(t)}^n$ has Hilbert polynomial $P$ for all $t\in T$}\right\}
    \end{align*}
    which we will call the \textbf{Hilbert scheme} once we know it is represented by a scheme projective over $S$.
\end{defn}



\section{Some Theory}

We want to think of all contravariant functors $h: (\Sch/S)^\op \to \Set$ as spaces, like $h_X = \Hom(-,X)$. One immediate necessary condition for such a contravariant functor to be representable is it being a \textit{Zariski sheaf}.

\begin{defn}
    A functor $h:()\Sch/S)^\op \to Set$ is called a \textbf{Zariski sheaf} if given an open cover
    \[X = \bigcup_i U_i,\]
    we get an equalizer sequence
    \[\bullet \to h(X) \to \prod_i h(U_i) \rightrightarrows \prod_{i,j} h(U_i\cap U_j).\]
\end{defn}

We will then carry over properties originally defined for morphisms of schemes to natural transformations $h'\to h$ by saying $h'\to h$ has property $P$ whenever base changed along a scheme $X$, we get a morphism of schemes with property $P$.

That is, given a natural transformation $h_X \to h$ (where $h_x = \Hom(-,X)$), we can then form the pullback
\[\begin{tikzcd}
h_X \times_h h' \dar\rar & h_X \dar\\
h' \rar & h
\end{tikzcd}\]
(where the pullback of functors is by pulling back in the codomain category, $\Set$, in each object), and then we ask that $h_X\times_h h'$ is representable by a scheme $Y$, and the morphism $h_Y\to h_X$ corresponds to a morphism
\[Y\to X\]
with property $P$.

\begin{defn}
    \begin{itemize}
        \item For two functors $h',h: (\Sch/S)^\op \to \Set$, we say a natural transformation $h' \to h$ expresses $h'$ as a \textbf{subfunctor} of $h$ if each component $h'(X)\to h(X)$ is injective.
        
        \item We say $h' \hookrightarrow h$ is an \textbf{open subfunctor} if it base changes to open immersions, i.e., given $h_X\to h$, $h_X \times_h h' \cong h_U$ for some scheme $U$, and the morphism $h_U\to h_X$ is an open immersion.
        
        \item We say a collection $h_i$ of open subfunctors \textbf{covers} $h$ if it base changes to open covers, i.e., given $h_X\to h$, for the $U_i\hookrightarrow X$ by
        \[h_{U_i} \cong h_X \times_h h_i \hookrightarrow h,\]
        the $U_i$ form an open cover of $X$.
    \end{itemize}
\end{defn}

\begin{prop}
    A locally representable (having an open cover by representable functors) Zariski sheaf is representable.
\end{prop}
\begin{proof}
    Glue the open subfunctors $h_{U_i}$ along their inclusion into $h$, i.e., we can define $h_{U_{ij}}$ as the fiber product
    \[h_{U_i} \times_h h_{U_j}\]
    and glue the $U_i$ along the $U_{ij}$. This gives us a scheme $X$.

    Next, build a morphism $\alpha: h_X\to h$, by seeing that this is the data (by Yoneda) of an $\alpha \in h(X)$, which we get exactly by the Zariski sheaf property, since $\alpha$ will push to the tuple of inclusions $\iota_i \in h(U_i)$, which are cooked up to agree on overlaps.

    Finally, check that $\alpha$ is an isomorphism. From its construction, it isn't hard to see that it is an isomorphism when base changed to each $U_i$. Then, to check when base changing to an arbitrary $T$, {\color{red} finish this!!}
    
    by seeing injective and surjective. Pick a test scheme $T$, so that we will study the set function $h_X(T) \to h(T)$.

    First, show surjective. Pick $\beta \in h(T)$ to try to reach. This is (by Yoneda) the data of $\beta: h_T \to h$. Base change this along $X$, and we want to argue that $h_X \times_h h_T \cong h_T$ via the projection, so that our pullback diagram is
    \[\begin{tikzcd}
    h_T \dar["="] \rar & h_X \dar\\
    h_T \rar & h
    \end{tikzcd}\]
    which factors $h_T\to h$ through $h_X$, providing a preimage of $\beta$ in $h_X(T)$.

    So, to see $h_X\times_h h_T \cong h_X$, we che
    
    To do this, recognize that by Yoneda, $\iota: h_U \to h$ is the data of an element $\iota \in h(U)$, and then 
\end{proof}

This then gives us a strategy to show certain moduli functors are representable---find an open cover (after you prove it is a Zariski sheaf).


\section{Grassmannian}

TODO: prove the Grassmannian is projective by Plucker embedding. Do the relative version. Also prove projective by proving proper and Chow's lemma lemma???


% \nocite{*} % Include all references in refs.bib regardless of whether they are referenced or not
% \bibliographystyle{plain} % We choose the "plain" reference style
% \bibliography{hilbert_scheme} % Entries are in the refs.bib file

\end{document}
